\section{Conclusions}
\label{ch:conclusions}

Software architecture recovery and dependency analysis represent fundamental
pillars of modern software engineering, underpinning critical activities
ranging from architectural smell detection and refactoring recommendation to
modularity assessment and technical debt monitoring. However, the accelerating
adoption of multi-language ecosystems and polyglot monorepos has exposed a
critical vulnerability in traditional static analysis tools: their heavy
reliance on compiler-coupled frontends and complex build-system bootstrapping.
In enterprise environments and large-scale empirical studies, requiring
fully resolved classpaths, compiler SDK installations, and synchronized build
descriptors leads to prohibitive operational failure rates, brittle tooling,
and severe scalability bottlenecks.

To address these limitations, this dissertation designed, formalized, and
empirically validated \textbf{OmniDeps}, a high-performance, language-agnostic
static analysis framework that decouples architectural dependency extraction
from native compiler toolchains. Operating under the foundational premise
that the analyzed source code compiles successfully within its native build
environment, OmniDeps demonstrates that Concrete Syntax Trees (CSTs) can be
directly queried and projected into normalized architectural dependency graphs
without executing compiler type-checking passes or whole-program bytecode
verification.

This research delivered three primary contributions:
\begin{enumerate}
  \item \textbf{Formalization of a Universal Intermediate Representation}:
    We established a strongly typed, mathematically grounded Intermediate
    Representation ($\mathcal{D}_{IR}$) based on algebraic data types, coupled
    with an algebraic query language governed by three orthogonal primitives
    (\texttt{Find}, \texttt{Extract}, and \texttt{Call}). This formalization
    abstracts structural declarations, modular scoping, and behavioral
    invocations away from idiosyncratic language grammars.
  \item \textbf{A Deterministic, High-Throughput Functional Pipeline}:
    We engineered an arena-allocated, parallel analysis pipeline implemented
    in native Rust. Leveraging Tree-sitter's incremental C parsers and Rayon's
    work-stealing concurrency, the architecture achieves single-pass CST
    extraction and deterministic two-phase symbol resolution across five
    diverse programming languages: C, C++, Java, Python, and Rust.
  \item \textbf{Rigorous Empirical Validation and Graph Harmonization}:
    To overcome the fundamental challenge of comparing heterogeneous static
    analyzers, we formulated a mathematical Graph Alignment and Harmonization
    Framework based on canonical entity identification, semantic entity
    lifting ($\lambda$), intermediate package synthesis, and unified edge
    taxonomies. Using this framework, OmniDeps was validated against formal
    micro-benchmark oracles and benchmarked against three state-of-the-art
    analyzers (\textit{Arcan}, \textit{Depends}, and \textit{Entangle}) across
    an extensive corpus of \textbf{156 open-source repositories} comprising over
    45 million lines of code.
\end{enumerate}

The empirical evaluation provided decisive validation of the proposed model:
\begin{itemize}
  \item \textbf{Fidelity}: On hand-crafted semantic ground truth, OmniDeps
    achieved \textbf{100\% precision, 100\% recall, and an $F_1$-score of
    1.000} across all supported languages, verifying the formal correctness of
    its scope climbing and query evaluation algorithms.
  \item \textbf{Parity}: In structured object-oriented ecosystems (Java),
    OmniDeps recovered approximately \textbf{98\% of architectural entities}
    identified by full compiler symbol tables, with edge Jaccard similarity
    converging to $\approx 0.58$--$0.60$ across independent baseline tools.
  \item \textbf{Operational Feasibility}: While compiler-coupled tools suffered
    severe attrition due to missing build configurations (Arcan completing only
    54.0\% of projects), OmniDeps achieved an overall completion rate of
    \textbf{99.4\%} (155 out of 156 repositories), demonstrating radical
    practical robustness.
  \item \textbf{Scalability}: Sustaining an average throughput of \textbf{74,445
    lines of code per second}, OmniDeps yielded speedups of \textbf{13.2x over
    Depends}, \textbf{48.2x over Arcan}, and \textbf{4.0x over Entangle},
    enabling near-instantaneous architectural extraction on systems with
    hundreds of thousands of lines of code.
\end{itemize}

In conclusion, this work proves that language-agnostic architectural dependency
extraction is both theoretically sound and industrially viable. By trading
redundant compiler-level type verification for declarative CST pattern queries
and decoupled lexical scope trees, OmniDeps bridges the gap between polyglot
versatility, architectural fidelity, and real-time execution performance.
